%% kapa: exam=1 round=33 year=2022 period=1 form=A subject=real_estate subject_ko=부동산학원론 seq=15 number=95
%% source: Q-Net 감정평가사 기출문제 (cst00309) artlSeq=5209997 file=raw/1008/5209997/2022년도 제33회 감정평가사 1차시험 1교시(민법,경제학,부동산학).pdf
%% conversion: kordoc 4.24.51 -> tools/gen1.py (auto); review=proofread-sonnet; flags=none
\begin{question}{95}
부동산투자이론에 관한 설명으로 옳지 않은 것은?
\begin{choices}
  \item 변동계수는 수익률을 올리기 위해 감수하는 위험의 비율로 표준편차를 기대수익률로 나눈 값이다.
  \item 포트폴리오를 구성하면 비체계적 위험을 회피할 수 있다.
  \item 위험기피형 투자자는 위험부담에 대한 보상심리로 위험할증률을 요구수익률에 반영한다.
  \item 두 개별자산으로 구성된 포트폴리오에서 자산간 상관계수가 양수인 경우에 음수인 경우보다 포트폴리오 위험절감효과가 높다.
  \item 투자안의 기대수익률이 요구수익률보다 높으면 해당 투자안의 수요증가로 기대수익률이 낮아져 요구수익률에 수렴한다.
\end{choices}
\end{question}
\answer{4}
